% Section 10.1, from lectures/L10/appendix/b_review.md.

\section{Course Review}
\label{c10:sec:course}\label{c10:app:b}
Every layer built in this book, from Chapter~\ref{c1:ch:linreg}'s linear regression up through the
CNN layers of Chapters~\ref{c6:ch:cnn} to~\ref{c10:ch:flatten}, follows the same three-step
pattern: \textbf{feedforward} (compute a prediction from an input), \textbf{backpropagate} (compute
how wrong that prediction was, and route the blame backward), and \textbf{optimize} (adjust
trainable parameters to reduce that error next time). What changes from layer to layer is
\emph{what} the parameters are, \emph{how} the feedforward sum is computed, and \emph{whether} there
are trainable parameters at all.

\subsection{Every layer type, side by side}
\label{c10:sec:sidebyside}
The first table says what each layer type is and where the book builds it; the second what each of
its three methods does.

\begin{booktable}{@{}>{\raggedright\arraybackslash}p{9.6em}p{6.4em}L@{}}
  \thead{Layer} & \thead{Chapters} & \thead{Trainable parameters} \\
  \midrule
  \code{Fixed} (linear regression) & 1, 2 & weight $k$, bias $m$ \\
  \code{Dense} & 3 to 5 & weights and a bias per node \\
  \code{Conv} & 6 to 8 & one shared kernel and a bias \\
  \code{MaxPool} & 6, 7, 9 & none \\
  \code{Flatten} & 6, 7, 10 & none \\
\end{booktable}

\begin{booktable}{@{}p{4.2em}LLL@{}}
  \thead{Layer} & \thead{Feedforward} & \thead{Backpropagate} & \thead{Optimize} \\
  \midrule
  \code{Fixed} & \code{y = k*x + m} & \code{δ = y_ref - y_p} & \code{k,m += δ * LR * (x or 1)} \\
  \addlinespace
  \code{Dense} & \code{y = σ(b + Σ w*x)} & \code{Δe = δ * σ'(s)} per node &
    \code{w += Δe * LR * input}, \code{b += Δe * LR} \\
  \addlinespace
  \code{Conv} & same as \code{Dense}, but local (a small window) and shared (the same kernel at
    every position) & gradients accumulated across every position the kernel visited, then applied
    once & \code{kernel,bias += gradient * LR} \\
  \addlinespace
  \code{MaxPool} & passes the largest value in each block forward & routes the incoming gradient
    only to the position that held the max value; everywhere else gets \code{0} & no-op: nothing to
    update \\
  \addlinespace
  \code{Flatten} & reshapes 2D to 1D & reshapes 1D to 2D (the exact inverse) & doesn't exist in the
    interface: there's nothing to optimize \\
\end{booktable}

A few things worth noticing across the rows:
\begin{itemize}
  \item \textbf{Only \code{Conv} and \code{Dense} have trainable parameters.} \code{MaxPool} and
    \code{Flatten} are purely structural: they reshape or downsample data, but never learn anything
    themselves.
  \item \textbf{\code{Conv} is \code{Dense} with two constraints added}: locality (each output only
    looks at a small window of the input) and weight sharing (the same kernel is reused everywhere
    instead of a unique weight per input and output pair). That's the entire conceptual jump from
    Chapter~\ref{c5:ch:dense} to Chapters~\ref{c6:ch:cnn} and~\ref{c7:ch:layers}: see
    \secref{c6:sec:images} if you want the full ``why'' behind those two constraints.
  \item \textbf{\code{Fixed} is a \code{Dense} layer with one node, no activation function, and no
    hidden layer around it.} The same \emph{error, adjustment, update} shape from
    Chapter~\ref{c1:ch:linreg} reappears, unchanged in spirit, all the way through the CNN layers.
\end{itemize}

\subsection{The full pipeline}
\label{c10:sec:pipeline}
\code{ml::cnn::Cnn}, in \file{lectures/L10/cnn_demo}, the finished reference implementation, wires
four of these layer types together into one working classifier. This is the same pipeline you built
up yourself as \code{cnn_work} across Chapters~\ref{c8:ch:conv} to~\ref{c10:ch:flatten}; worth
comparing the two if anything in your own version doesn't behave the way you expect:

\begin{console}
input → Conv → MaxPool → Flatten → Dense → prediction
\end{console}

\noindent\textbf{Feedforward} runs left to right, each layer consuming the previous layer's
\code{output()}.

\textbf{Backpropagation} runs right to left: the target vector goes into the \code{backpropagate()}
of \code{Dense} first, and each earlier layer is then fed the next layer's
\code{inputGradients()}: \code{Dense → Flatten → MaxPool → Conv}.

\textbf{Optimization} updates every layer's parameters (\code{Conv} and \code{Dense} only) using the
gradients already computed by \code{backpropagate()}.

\file{lectures/L10/cnn_demo/README.md} has the full method-by-method breakdown, including the exact
formulas used in the C++ implementation; \secref{c8:sec:pipeline} is the same breakdown for
\code{cnn_work}.
